\begin{minipage}{0.75\linewidth}
    On réalise une expérience en plaçant un émetteur $E$ d'ultrasons à une
    distance $d$ d'un récepteur $R$ d'ultrasons. L'émetteur $E$ émet à l'instant
    $t_{0}=0$ un signal ultrasonore de fréquence $N=\qty{40}{\kHz}$, ce signal
    est reçu par $R$ après un retard $\tau$.
    \begin{enumerate}
        \item Les ondes ultrasonores sont-elles des ondes mécaniques~?
              Justifier.
        \item Pour différentes valeurs de $d$, on mesure le retard $\tau$. Le
              graphe de la figure~\ref{fig:tau_d}, donne la variation de $\tau$
              en fonction de $d$.
        
              En exploitant le graphe, déterminer la valeur de la célérité $v$
              des ondes ultrasonores.
        \item Déduire la valeur de la longueur d'onde $\lambda$ des ultrasons.
    \end{enumerate}
\end{minipage}
\hfill
\begin{minipage}{0.23\linewidth}
    \begin{figure}[H]
        \centering
        \includegraphics[width=\textwidth]{2025-11-01_015224-}
        \caption{}
        \label{fig:tau_d}
    \end{figure}
\end{minipage}




